% Forkcast — Economic Impact Evidence Brief (NIW support) — July 2026
% Compile: latexmk -pdf forkcast-niw-economic-impact.tex
\documentclass[10pt]{article}

\usepackage[letterpaper,margin=0.95in,top=0.9in,bottom=1in]{geometry}
\usepackage[utf8]{inputenc}
\usepackage[T1]{fontenc}
\usepackage{helvet}\renewcommand{\familydefault}{\sfdefault}
\usepackage{microtype}
\usepackage{xcolor}
\usepackage{booktabs}
\usepackage{tabularx}
\usepackage{enumitem}
\usepackage{titlesec}
\usepackage{fancyhdr}
\usepackage[hidelinks]{hyperref}

\definecolor{accent}{HTML}{EC3013}
\definecolor{ink}{HTML}{201E1D}
\definecolor{ground}{HTML}{F3F2F2}
\definecolor{muted}{HTML}{605D5D}

\color{ink}
\setlength{\parindent}{0pt}
\setlength{\parskip}{5pt}
\renewcommand{\arraystretch}{1.35}

\titleformat{\section}
  {\Large\bfseries\color{ink}}
  {\textcolor{accent}{\thesection}}{0.7em}{}
  [\vspace{-6pt}{\color{accent}\rule{2.2cm}{2pt}}]
\titlespacing*{\section}{0pt}{16pt}{8pt}
\titleformat{\subsection}{\normalsize\bfseries\color{ink}}{\thesubsection}{0.6em}{}
\titlespacing*{\subsection}{0pt}{9pt}{3pt}

\pagestyle{fancy}
\fancyhf{}
\fancyhead[L]{\small\color{muted}\textbf{\textcolor{accent}{Fork}cast} --- Economic Impact Evidence Brief}
\fancyhead[R]{\small\color{muted}July 2026 · For immigration counsel}
\fancyfoot[C]{\small\color{muted}\thepage}
\renewcommand{\headrulewidth}{0.4pt}
\renewcommand{\headrule}{{\color{ground}\hrule width\headwidth height 1pt}}

\newcolumntype{L}{>{\raggedright\arraybackslash}X}

\begin{document}

% ---- Header ----
{\color{accent}\rule{2.6cm}{5pt}}\\[10pt]
{\huge\bfseries Economic Impact Evidence Brief}\\[4pt]
{\large\color{muted}Supporting co-founder Seymur Hasanov's EB-2 NIW petition · July 2026}\\[8pt]

\begin{tabularx}{\textwidth}{@{}L@{}}
\toprule
\textbf{Framing.} Prepared by the founders as source material for immigration counsel. This document is not legal
advice and draws no legal conclusions; it organizes the factual, economic, and documentary record so qualified
counsel can assess and argue it under the applicable standard (\emph{Matter of Dhanasar}). Every market figure is
sourced (CDC, USDA ERS, AHRQ, peer-reviewed literature; full citations in the Market Research Dossier). Projections
are labeled as projections. Nothing herein claims users, partners, revenue, or health outcomes that do not exist.\\
\bottomrule
\end{tabularx}

% ================= 1 =================
\section{The proposed endeavor --- specific and defined}

To build and scale \textbf{Forkcast}, a U.S. nutrition decision-support and data-verification system for restaurant
meals, with three components:

\begin{enumerate}[leftmargin=1.4em,itemsep=2pt]
\item \textbf{A consumer product} that scores restaurant dishes against an individual's clinically-derived targets
\emph{before} ordering, then closes the loop: order, confirmation, evidence-carrying meal log, and adaptive
recalibration from the user's own measured energy balance.
\item \textbf{A small-business data program} that digitizes, estimates, and owner-verifies nutrition data for
\textbf{independent restaurants} --- the majority of U.S. restaurants, which federal menu-labeling rules
(21 CFR 101.11; chains of 20+ locations only) do not reach --- under a documented protocol with versioned public
corrections.
\item \textbf{A measurement framework}: six pilot metrics pre-registered publicly \emph{before} any measurement,
so national-scale claims are earned from reproducible evidence, not asserted.
\end{enumerate}

The endeavor is \textbf{operational today}: a complete working product with a dated 2026 development record,
three-tier data provenance including live-captured restaurant-published nutrition (Sweetgreen; The Halal Guys), a
restaurant-side order terminal with versioned corrections, and role-based accounts for diners and restaurants.

% ================= 2 =================
\section{The national problem, in economic terms}

\begin{tabularx}{\textwidth}{@{}L r >{\raggedright\arraybackslash}p{4.2cm}@{}}
\toprule
\textbf{Economic fact} & \textbf{Figure} & \textbf{Source}\\
\midrule
Annual U.S. cost of diet-related disease & \textbf{$\sim$\$334B/yr} & CDC; AHRQ MEPS\\
--- of which obesity-attributable medical cost & \textbf{$\sim$\$173B/yr} (2019\$) & CDC\\
U.S. adults with obesity / overweight & \textbf{40.3\% + 31.7\%} $\approx$ 72\% & CDC NCHS, 2024\\
Food spending away from home & \textbf{58.5\%} (record; \$4,485/capita) & USDA ERS, 2024\\
Daily calories from food away from home & $\sim$1/3 & USDA ERS / NHANES\\
Sodium density, restaurant vs.\ home food & 2,151 vs 1,369 mg/1,000 kcal & USDA ERS EIB-105\\
Effect of federal menu calorie labels & \textbf{$\sim$22--25 kcal/order}; often null & Peer-reviewed literature\\
Consumers underestimating restaurant calories & \textbf{2 in 3}; $\sim$25\% by 500+ kcal & USDA/NHANES 2024\\
\bottomrule
\end{tabularx}

\vspace{4pt}
\textbf{The economic syllogism:} more than half of U.S. food spending flows through a channel that measurably worsens
diets, largely escapes the federal transparency requirement, and has resisted the one intervention tried at scale
(passive labels). The cost of that failure --- $\sim$\$334B a year --- is national by any definition. This endeavor
operates \emph{upstream of the order}, covers the \emph{unregulated independent majority}, and \emph{measures its own
effect} --- the three points where existing policy and products do not.

% ================= 3 =================
\section{Channels of economic impact}

\subsection*{3.1 Public-health cost reduction (primary channel, honestly framed)}
The mechanism targets the decision moment where labels fail: personalized, explained scoring before ordering; allergy
and condition advisories; and a confirmed-order log that makes dietary change measurable per user. \textbf{No
health-outcome effect is claimed today.} The pre-registered \emph{diet-quality delta} metric (change in sodium and
calories per restaurant meal vs.\ each user's own baseline) is designed to quantify the effect. Context for scale:
labeling policy's $\sim$24 kcal/transaction justified federal rulemaking; a decision-support tool need only exceed
that low bar, per user, to outperform the status-quo intervention --- and the framework will show whether it does.

\subsection*{3.2 Small-business competitiveness (independent restaurants)}
\begin{itemize}[leftmargin=1.3em,itemsep=2pt]
\item \textbf{Free digitization and verification tools} (order terminal, allergy flags, versioned menu corrections)
deliver capabilities chains buy from enterprise vendors --- at no cost to the independent majority.
\item \textbf{A cheaper demand channel:} 6\% commission only on orders the platform originates, vs.\ the 15--30\%
independents pay delivery marketplaces (the market's own pickup benchmark is 6\%) --- direct margin relief.
\item \textbf{Portable reputation infrastructure:} verified nutrition data and public correction history for small
businesses that otherwise have none.
\item \textbf{Priced, replicable onboarding:} the pilot measures the cost of digitizing one independent restaurant
--- a pre-registered metric.
\end{itemize}

\subsection*{3.3 Direct economic activity (projections, labeled)}
From the planning model (base case, July 2026): revenue $\sim$\$0.2M (Y1) $\to$ $\sim$\$37M ARR (Y5); employment
growing from $\sim$5--6 to an estimated \textbf{$\sim$45--60 U.S. jobs} by Year 5 (engineering, dietetics, operations,
restaurant partnerships), plus indirect activity at partner restaurants; payroll and corporate tax base, and state
meals-tax collections on platform-originated orders (the product computes Massachusetts meals tax on every order
today). The protocol is market-agnostic by design --- Boston is the first instance, not the boundary.

\subsection*{3.4 Non-dilutive R\&D alignment}
The estimation-engine research agenda maps to standing federal SBIR priorities (NIH/NIDDK metabolic-disease
prevention; USDA-NIFA food science \& nutrition; NSF). Submissions are planned, not claimed.

\subsection*{3.5 Knowledge and infrastructure spillovers}
A publishable menu-collection protocol; an order-ground-truth validation dataset for nutrition-estimation accuracy
that retrospective logging products structurally cannot build; and a transferable correction-governance pattern
(versioned, timestamped, public) for consumer-data integrity in AI products.

% ================= 4 =================
\section{Factual mapping to the Dhanasar framework}
{\small\color{muted}Counsel determines legal sufficiency; this table organizes facts under each prong.}

\begin{tabularx}{\textwidth}{@{}>{\raggedright\arraybackslash}p{4.0cm} L@{}}
\toprule
\textbf{Prong} & \textbf{Factual support assembled}\\
\midrule
1. Substantial merit \& national importance & Sourced economics of \S2; the federal transparency gap for independent
restaurants; documented failure of passive labeling; impact channels of \S3; a protocol designed for cross-market
U.S. scale rather than one locality.\\
2. Well positioned to advance the endeavor & Petitioner (co-founder, Harvard SEAS engineering); hands-on builder of a complete working
product (dated 2026 record); demonstrated data capability (live published-nutrition integrations, three-tier
provenance, adaptive energy-balance calibration); documented methods and verified-figure business plan; pilot
designed and publicly pre-registered; funding strategy mapped (pre-seed + SBIR).\\
3. Balance favoring waiver \emph{(as argued by counsel)} & Founder-created enterprise, co-founded by the petitioner (job-creating); work product ---
protocol, dataset, correction governance --- is infrastructure others can build on; urgency: diet-disease costs
compound annually while the independent-restaurant gap persists.\\
\bottomrule
\end{tabularx}

% ================= 5 =================
\section{Evidence inventory}

\subsection*{Exists today (dated, verifiable)}
Working end-to-end product with 45+ dated commits (2026) and a public demo deployment · three-tier provenance catalog
with live-captured published nutrition (retrieval dates and non-affiliation statements on each listing) · public
pre-registered pilot metrics and source ledger · documented methods (menu-collection protocol; system architecture)
· business plan v2 with all 25 reused market figures verified against the sourced dossier; financial model v2 ·
competitive scan with dated live captures · founders' delivery-cost research validated with restaurant owners ·
product-embedded integrity system (prototype stamps, simulation labels, confidence tiers, no-unrecorded-claims
commitment).

\subsection*{To be generated by the pilot (the strengthening path)}
Signed verification letters from 3--5 Boston independent restaurants with their correction logs · measured pilot
metrics against the pre-registered definitions · expert letters (registered dietitian; health economist;
independent-restaurant operator or association; public-health academic) · SBIR submission and, if awarded, the award
· the order-ground-truth validation dataset and any resulting write-up · pilot-grounded press coverage.

% ================= 6 =================
\section{Stated weaknesses (before anyone else states them)}

\begin{itemize}[leftmargin=1.3em,itemsep=2pt]
\item \textbf{No measured outcomes yet} --- mitigated by pre-registration: definitions were published before data.
\item \textbf{Lean founding-team execution risk} --- mitigated by the shipped record of the two-person founding team; hiring is specified and funded.
\item \textbf{Projections are projections} --- every forward number is labeled; the base case was deliberately
\emph{reduced} for defensibility (\$58M $\to$ \$37M), itself evidence of conservatism.
\item \textbf{Catalog's real restaurants are not partners} --- stated on every listing; converting them is the
pilot's first milestone, with letters as the artifact.
\end{itemize}

\vspace{10pt}
{\color{ground}\hrule height 2pt}
\vspace{4pt}
{\small\color{muted}Companions: Business Plan v2 · Financial Model v2 · Market Research Dossier (citations for every
figure above) · Menu Collection Protocol · Competitive Scan (July 2026). To be reviewed, supplemented, and framed by
qualified immigration counsel before any filing.}

\end{document}
